\section{Complex-zero motion and transport from a collision}
\label{sec:complex-dynamics}

The local coefficients become useful for statistics only after one knows
where a nonreal pair reaches the zeta slice. We therefore separate the
exact motion laws from the estimates needed to continue them. Throughout
this section, a finite cluster is isolated by a zero-free boundary;
simple-root equations are used only on intervals of simple labeling.

\subsection{Finite-cluster dynamics}
Factor \(H_\lambda=P_\lambda U_\lambda\), with
\(P_\lambda(z)=\prod_j(z-z_j(\lambda))\) the monic cluster polynomial
and \(U_\lambda\) nonvanishing in its neighborhood. Put
\(g=\partial_z\log U_\lambda\). At a simple root, differentiating
\(H_\lambda(z_j(\lambda))=0\) and using
\(\partial_\lambda H=-H_{zz}\) gives
\begin{equation}\label{eq:complex-root-motion}
 z_j'=2\sum_{k\ne j}\frac1{z_j-z_k}+2g(z_j),
 \qquad {}'=\partial_\lambda.
\end{equation}
Indeed \(H_{zz}/H_z=P_{zz}/P_z+2U_z/U\) at the root, and
\(P_{zz}(z_j)/P_z(z_j)=2\sum_{k\ne j}(z_j-z_k)^{-1}\).
Thus no infinite zero sum is needed for this identity. Writing
\(z_j=x_j+\ii y_j\) yields
\begin{align}
 x_j'&=2\sum_{k\ne j}\frac{x_j-x_k}{|z_j-z_k|^2}
                              +2\Re g(z_j),\notag\\
 y_j'&=-2\sum_{k\ne j}\frac{y_j-y_k}{|z_j-z_k|^2}
                              +2\Im g(z_j).
 \label{eq:planar-root-motion}
\end{align}
The transverse interaction has the opposite sign. In particular,
the real inverse-square interaction cannot remain a positive energy
after complexifying the gaps: for a conjugate pair it is
\((2\ii y)^{-2}=-1/(4y^2)\). An absolute-square replacement defines
a different quantity and needs its own evolution argument.

\begin{proposition}[Transverse balance]\label{prop:transverse-balance}
On a simple-root interval for the isolated cluster,
\begin{equation}\label{eq:transverse-balance}
 \frac{d}{d\lambda}\sum_j y_j^2
 =-4\sum_{j<k}\frac{(y_j-y_k)^2}{|z_j-z_k|^2}
                      +4\sum_j y_j\Im g(z_j).
\end{equation}
For a conjugate pair \(x\pm\ii y\), \(y>0\),
\begin{equation}\label{eq:conjugate-pair-height}
 (y^2)'=-2+4y\Im g(x+\ii y).
\end{equation}
\end{proposition}
\begin{proof}
Multiply the second line of \eqref{eq:planar-root-motion} by \(2y_j\)
and combine the two ordered terms for each pair. Their numerator is
\(y_j(y_j-y_k)+y_k(y_k-y_j)=(y_j-y_k)^2\).
For two conjugate roots the mutual term in \eqref{eq:complex-root-motion}
is \(-\ii/y\); hence \(y'=-1/y+2\Im g\).
\end{proof}
For a real polynomial heat flow containing all its roots, \(g=0\),
so this transverse second moment decreases in Newman time. For a local
cluster, the exterior term has no prescribed sign.

\subsection{An exact pair equation and conditional tube bounds}
Let \(\lambda_*>0\) be a collision time and set
\(s=\lambda_*-\lambda\). Increasing \(s\) is forward heat time;
the zeta slice is reached at \(s=\lambda_*\). Suppose an isolating
neighborhood persists on \(0\le s\le S\) and contains exactly the
pair born at the collision. In that neighborhood write
\begin{equation}\label{eq:tube-factor}
 H_{\lambda_*-s}(z)=[(z-x(s))^2+y(s)^2]U(s,z),\qquad
 g(s,z)=U_z/U,
\end{equation}
with \(U\) real on the real segment and nonzero. The roots are assumed
simple with \(y(s)>0\) for \(s>0\). From
\eqref{eq:complex-root-motion},
\begin{equation}\label{eq:pair-s-motion}
 x_s=-2\Re g(s,x+\ii y),\qquad
 (y^2)_s=2-4y\Im g(s,x+\ii y).
\end{equation}
The singular mutual force cancels from the midpoint equation.

Assume the full vertical segment from \(x\) to \(x+\ii y\) lies in
the nonvanishing region of \(U\). Since \(g(s,x)\) is real,
\begin{equation}\label{eq:tube-k}
 k(s):=\frac{\Im g(s,x+\ii y)}{y}
       =\int_0^1\Re g_z(s,x+\ii r y)\dd r.
\end{equation}
Putting \(u=y^2\), the transverse equation becomes the scalar linear
equation \(u_s+4k(s)u=2\), with \(u(0)=0\). Its exact solution is
\begin{equation}\label{eq:tube-integrating-factor}
 u(s)=2\int_0^s\exp\!\left(-4\int_r^s k(q)\dd q\right)\dd r.
\end{equation}
This identity expresses the entire environmental correction in one
integrated coefficient.

\begin{theorem}[Conditional pair transport]\label{thm:tube}
Under the preceding isolation and segment hypotheses, suppose
\(|g_z|\le K\) on those segments and
\(|\Re g(s,x+\ii y)|\le V\) along the pair, for all \(0\le s\le S\).
Then
\begin{equation}\label{eq:tube-bounds}
 \frac{1-\ee^{-4Ks}}{2K}\le y(s)^2
 \le\frac{\ee^{4Ks}-1}{2K},\qquad
 |x(s)-x_*|\le2Vs.
\end{equation}
When \(K=0\), both transverse expressions mean \(2s\).
If the tube reaches \(S=\lambda_*\), the final pair has transverse
size at least
\begin{equation}\label{eq:final-tube-data}
 d_*:=\left(\frac{1-\ee^{-4K\lambda_*}}{2K}\right)^{1/2},
 \qquad \left|\theta-\frac{x_*}{2}\right|\le V\lambda_*,
\end{equation}
where \(\theta=x(\lambda_*)/2\) is its zeta ordinate.
\end{theorem}
\begin{proof}
Equation~\eqref{eq:tube-k} gives \(|k|\le K\). Bound the exponential
in \eqref{eq:tube-integrating-factor} between
\(\ee^{-4K(s-r)}\) and \(\ee^{4K(s-r)}\) and integrate.
Integrating the midpoint equation proves the remaining inequalities.
The relation between \(\theta\) and the midpoint follows from
\eqref{eq:zeta-H-zero-map} below.
\end{proof}

At the event \(g=\ell\) and \(g_z=\kappa\), using
\eqref{eq:ell-kappa}. Analyticity of the symmetric pair coordinates
and \eqref{eq:tube-k} give \(k(s)=\kappa+O(s)\), whence
\begin{equation}\label{eq:postcollision-jet}
 x(s)=x_*-2\ell s+O(s^2),\qquad
 y(s)^2=2s-4\kappa s^2+O(s^3).
\end{equation}
For the second formula, substitute \(u=2s+cs^2+O(s^3)\) into
\(u_s=2-4ku\), obtaining \(c=-4\kappa\).
The midpoint velocity \(-2\ell\) differs from the critical-point
velocity \(-\beta=-3\ell\). Interchanging them would give an
incorrect final-height prediction.

The theorem is a transport estimate on an already controlled tube,
not a proof that a local event neighborhood reaches time zero.
Its use requires a zero count of two, a nonzero boundary margin,
complete-source errors and uniform bounds for \(g,g_z\) along the
whole proposed path. The contour interface in
Section~\ref{sec:coefficients} supplies a way to estimate these
quantities. Keeping the resulting enclosures strictly inside the
controlled region is the additional continuation requirement.

\subsection{Strip width and a conditional environmental amplification}
For the canonical family, the full paired zero-dynamics formula is
also available in the complex regime
\cite[Proposition~3.1]{Polymath2019}; pairing opposite zeros fixes
the convergence convention. De Bruijn's strip theorem, in
\cite[Theorem~3.2]{Polymath2019}, gives, for the supremal imaginary
width \(Y(\lambda)\) and \(\lambda_2>\lambda_1\),
\begin{equation}\label{eq:strip-contraction}
 Y(\lambda_2)^2\le
 \max\{Y(\lambda_1)^2-2(\lambda_2-\lambda_1),0\}.
\end{equation}
The supremum version follows by decreasing an admissible strip bound.
A collision at \(\lambda_*>0\) implies
\(\Lambda\ge\lambda_*\), since nonreal roots occur immediately below
it. The same theorem gives \(\Lambda\le Y(0)^2/2\), hence
\begin{equation}\label{eq:global-defect}
 Y(0)\ge\sqrt{2\lambda_*}.
\end{equation}
This familiar consequence identifies some transverse defect. It neither
locates that defect nor follows the same pair to time zero; a supremum
need not be attained.

\begin{proposition}[Persistent real neighbors]\label{prop:neighbor-growth}
Suppose a simple upper root \(x(s)+\ii y(s)\), born at a real collision,
remains of maximal imaginary part among all zeros for \(0<s\le S\).
Assume that throughout this interval at least \(m\ge1\) real zeros,
counted with multiplicity, remain within horizontal distance \(R>0\)
of \(x(s)\). Then, with \(u=y^2\),
\begin{equation}\label{eq:neighbor-ode}
 u_s\ge2+\frac{4mu}{R^2+u}.
\end{equation}
While \(u\le R^2\),
\begin{equation}\label{eq:neighbor-exp}
 u(s)\ge\frac{R^2}{m}\bigl(\ee^{2ms/R^2}-1\bigr).
\end{equation}
If the hypotheses persist long enough, \(u\) reaches \(R^2\) no
later than \(s_R=R^2\log(1+m)/(2m)\); thereafter \(u_s\ge2+2m\).
\end{proposition}
\begin{proof}
Reverse the Newman-time full dynamics. Maximality makes each transverse
force \(2(y-y_k)/|z-z_k|^2\) nonnegative. Its conjugate partner
contributes \(1/y\) to \(y_s\), and each indicated real zero contributes
at least \(2y/(R^2+u)\). Multiplication by \(2y\) gives
\eqref{eq:neighbor-ode}. The imaginary-part series is absolutely
convergent: zeros lie in a strip and their distant contributions are
bounded by a constant times \(|z_k|^{-2}\). Thus retaining these
nonnegative terms is legitimate. For \(u\le R^2\),
\(u_s\ge2+(2m/R^2)u\); an integrating factor proves
\eqref{eq:neighbor-exp}. Monotonicity and
\eqref{eq:neighbor-ode} prove the hitting-time and subsequent bounds.
\end{proof}

If \(0<R\le1\) and these hypotheses persist to \(S=\lambda_*\),
the canonical open time-zero strip \(y^2<1\) necessarily implies
\begin{equation}\label{eq:neighbor-strip-budget}
 \lambda_*<\frac{R^2\log(1+m)}{2m}
                   +\frac{1-R^2}{2+2m}.
\end{equation}
Indeed the right-hand side is an upper bound on the time needed to
reach \(u=1\); equality would already violate the open strip.
This is a conditional comparison, not a numerical Newman bound.
Maximality may switch between roots, a complex collision may interrupt
simple labeling, and the nearby real population must persist.
Maximality starting at \(y=0\) also requires all other zeros to be
real at that event. None of these conditions follows from the local
normal form. Existing spatial barriers provide a rigorous alternative
\cite[Proposition~3.3]{Polymath2019}, but their initial, far-field and
intermediate-wall hypotheses must all be retained.

\subsection{The cost of a moving spatial window}
Localization introduces further terms even before any arithmetic
estimate is attempted. Let \(a_j(s)\) be differentiable real weights
on a fixed isolated finite cluster of simple roots, and use
\(s=\lambda_*-\lambda\).
\begin{proposition}[Weighted transverse balance]\label{prop:weighted-balance}
With the exterior factor of that cluster,
\begin{align}
 \frac d{ds}\sum_j a_jy_j^2
 ={}&4\sum_{j<k}
 \frac{(a_jy_j-a_ky_k)(y_j-y_k)}{|z_j-z_k|^2}
 -4\sum_j a_jy_j\Im g(z_j)
 +\sum_j (a_j)_s y_j^2.
 \label{eq:weighted-balance}
\end{align}
\end{proposition}
\begin{proof}
Reverse the sign of \eqref{eq:planar-root-motion}, multiply its
imaginary component by \(2a_jy_j\), and pair the \(j,k\) terms.
Differentiating the weights gives the last sum.
\end{proof}
For \(a_j=w((x_j-X(s))/R(s))\), with \(R(s)>0\),
\begin{equation}\label{eq:window-velocity}
 (a_j)_s=w'\!\left(\frac{x_j-X}{R}\right)
 \left(\frac{(x_j)_s-X_s}{R}-\frac{(x_j-X)R_s}{R^2}\right).
\end{equation}
Even nonnegative weights do not make the internal term nonnegative:
\begin{align*}
 (a_jy_j-a_ky_k)(y_j-y_k)
 =\frac{a_j+a_k}{2}(y_j-y_k)^2
       +\frac{a_j-a_k}{2}(y_j^2-y_k^2).
\end{align*}
The unequal-weight correction, window derivative and exterior force
each require an estimate. Following one pair with the center keeps its
own plateau weight fixed, but removes none of these costs for neighbors.
A path visiting many heights at different times also need not produce
a broad cluster on the final slice. Transport and final spatial
amplification are separate assertions.
